\documentclass{sdjummidterm}
\begin{document}
\begin{titlepage}
\centering
\vspace*{\dimexpr2mm-12bp\relax}
\CoverLogo
\vspace{12mm}
\CoverTitle{硕士研究生中期考核表}
\vspace{23mm}
\begin{coverfields}
\CoverField{学\quad 号}{local:midterm}
\CoverField{姓\quad 名}{测试研究生甲}
\CoverField{校内导师}{测试导师甲}
\CoverField{企业导师}{测试导师乙}
\CoverField{研究方向}{模拟设备状态监测}
\CoverField{类别（领域）}{电子信息（控制工程）}
\CoverField{学\quad 院}{测试工程学院}
\CoverField{考核时间}{2026年11月28日}
\end{coverfields}
\vfill
{\sdjufont{15}{23}上海电机学院研究生处\par 2026年11月28日\par}
\vspace{22mm}
\end{titlepage}

\FormHeading{一、自我小结}
\begin{formbox}{113mm}[思想品德、行为表现、综合素质]
\noindent（一）思想品德、行为表现、综合素质等\par
按模拟课题要求记录实验过程，参加组内复核。\par\vfill
\noindent 自我评价（优、良、合格、不合格）：良。
\end{formbox}
\begin{formbox}{92mm}[科研与学术活动]
\noindent（二）参加科研项目、学术会议、科创竞赛、专利申请情况，发表论文篇数\par
参加虚构校企项目1项，模拟学术交流1次。科创竞赛、专利申请及论文发表均为0。
\end{formbox}

\clearpage
\begin{formbox}{68mm}[奖惩情况]
\noindent（三）奖惩情况\par
无。
\end{formbox}
\FormHeading{（四）已修课程及成绩}
\begingroup
\setlength\tabcolsep{3bp}
\sdjufont{10.5}{16}
\begin{longtable}{|C{7mm}|C{34mm}|C{10mm}|C{10mm}|C{8mm}|C{34mm}|C{10mm}|C{10mm}|}
\hline
\multicolumn{4}{|c|}{学位课} & \multicolumn{4}{c|}{非学位课}\\\hline
 & 课程名称 & 学分 & 成绩 & & 课程名称 & 学分 & 成绩\\\hline
\endhead
\multicolumn{4}{|c|}{须填写所有已出成绩不得遗漏} & \multicolumn{4}{c|}{须填写所有已出成绩不得遗漏}\\\hline
\multirow{4}{*}{\VerticalLabel{学\\位\\课}} & 测试控制理论 & 3 & 90 & \multirow{4}{*}{\VerticalLabel{非\\学\\位\\课}} & 测试学术写作 & 2 & 92 \RowHeight{12mm}\\\cline{2-4}\cline{6-8}
 & 测试数据分析 & 3 & 93 & & 测试实验规范 & 1 & 95 \RowHeight{12mm}\\\cline{2-4}\cline{6-8}
 & & & & & & & \RowHeight{12mm}\\\cline{2-4}\cline{6-8}
 & & & & & & & \RowHeight{12mm}\\\hline
\multicolumn{8}{|l|}{学院确认：\Signature{dean}\RowHeight{15mm}}\\\hline
\end{longtable}
\endgroup

\clearpage
\FormHeading{二、学位论文}
\begin{formbox}{9mm}
\noindent 学位论文题目：\TitleText{面向模拟台架的设备运行状态监测与异常检测方法研究}
\end{formbox}
\begin{formbox}{79mm}[学位论文研究情况]
\noindent（一）学位论文研究情况和已完成的研究内容\par
完成模拟台架设计、工况定义和首批样本核验。
\end{formbox}
\begin{formbox}{50mm}[阶段性成果]
\noindent（二）所取得的阶段性成果\par
形成模拟数据字典、测试日志及基线结果。
\end{formbox}
\begin{formbox}{59mm}[开题内容对照]
\noindent（三）已完成的工作是否存在与开题报告内容不相符的部分，如存在请说明其原因\par
延后一组扩展试验，用于统一工况标记规则。
\end{formbox}

\clearpage
\begin{formbox}{99mm}[下一步工作计划]
\noindent（四）下一步工作计划和研究内容\par
补齐扩展工况，复核各批次测试配置和统计结果。
\end{formbox}
\begin{formbox}{104mm}[学术论文发表情况及计划]
\noindent（五）学术论文发表情况及拟发表论文内容和计划\par
当前无已发表论文。拟整理模拟台架异常检测的复现记录。
\end{formbox}

\clearpage
\FormHeading{三、考核意见}
\begin{formbox}{46mm}[导师审核意见][导\\师\\审\\核\\意\\见]
\noindent （是否同意参加中期考核答辩，对研究生的思想品德、课程学习、论文及科研情况进行综合评价）\par
同意参加中期考核答辩。\par\vfill
\begin{flushright}
签字：\Signature{supervisor}\\\SignatureDate{supervisor}
\end{flushright}
\end{formbox}
\begin{formbox}{40mm}[学院审核意见][学\\院\\审\\核\\意\\见]
\noindent （包括思想品德和学术规范遵守情况等）\par
同意提交考核小组。\par\vfill
\begin{flushright}
签字：\Signature{dean}\\\SignatureDate{dean}
\end{flushright}
\end{formbox}
\begin{formbox}{39mm}[考核小组意见][考\\核\\小\\组\\意\\见]
\noindent （包括论文及科研情况、综合素质等）\par
建议补充说明样本筛选规则。\par\vfill
\begin{flushright}
签字：\Signature{committee_chair}\\\SignatureDate{committee_chair}
\end{flushright}
\end{formbox}
\begin{formbox}{13mm}
\noindent 考核结果（请打“√”）\qquad\FormCheckBox\,合格\qquad\FormCheckBox\,不合格
\end{formbox}
\begin{longtable}{|C{14mm}|C{22mm}|C{17mm}|C{44mm}|C{26mm}|}
\hline
\multicolumn{5}{|c|}{考 核 小 组}\\\hline
 & 姓 名 & 职 称 & 单 位 & 签 字 \\\hline
\endhead
组长 & 测试专家甲 & 教授 & 测试工程学院 & \Signature[23mm]{committee_chair}\RowHeight{11mm}\\\hline
成员 & 测试专家乙 & 工程师 & 测试实践中心 & \Signature[23mm]{committee_member-1}\RowHeight{11mm}\\\hline
秘书 & 测试秘书甲 & 讲师 & 测试工程学院 & \Signature[23mm]{recorder}\RowHeight{11mm}\\\hline
\end{longtable}
\end{document}
